\section {Architecture}

The architecture of the Data Structures and Algorithms (DSA) Visualizer project is structured hierarchically to enforce a strict separation of concerns among the core application loop, user interface states, graphical rendering, and the underlying mathematical data models. At the highest entry and management level, the execution delegates the primary event loop, comprising state updates and rendering, to an application core. This core relies on a scene manager that governs the overarching application state by maintaining a collection of scenes via unique smart pointers, facilitating dynamic transitions between discrete functional areas of the software. Concurrently, a resource manager employs the Singleton design pattern to centralize the memory management and distribution of graphical assets, such as fonts and textures.\\

To decouple user input from system execution, the architecture implements a robust Command Pattern layer. This mechanism revolves around an abstract interface defining a pure virtual execution method, which is subsequently realized by specialized command implementations. Scene commands leverage factories and lambda functions to instantiate new user interface states dynamically, while operation commands encapsulate specific algorithmic manipulations requested within the visualizer interfaces.\\

The graphical interface is managed by a distinct scene layer anchored by a foundational base interface. This base dictates a standard lifecycle requiring derived classes to process events, update logical states, and render graphical components. This hierarchy branches into navigational elements, such as menu scenes containing interactive callback-driven buttons, and complex visualization scenes. These abstract visualization scenes serve as templates that house shared user interface components and define abstract contracts for algorithmic operations like insertion, searching, deletion, and traversal. Concrete implementations for specific paradigms---such as linked lists, graphs, trees, and hash sets---inherit from this abstract scene to implement the distinct user interactions required for each structural type.\\

Bridging the interactive interfaces and the raw algorithms is the data model and visualizer layer. A base visualizer interface establishes the mandatory protocols for resetting state and pushing graphical representations to the rendering window. Specific visualizer classes implement these protocols to translate raw data into dynamic on-screen components. At the foundational base of the entire architecture lie the pure data model files. These components encapsulate the intrinsic algorithmic logic of the data structures completely independently of the Simple and Fast Multimedia Library (SFML) or any graphical dependencies. The visualizers maintain references to these pure models, ensuring that the visual representation merely observes and reflects the state of the isolated algorithmic backend.\\

Given this architecture, it is worth considering how the system manages the temporal synchronization between the pure, instantaneous C++ data models and the sequential, frame-by-frame animations required by the visualizer layer. Furthermore, one might examine how the Command Pattern specifically handles the undoing of state-altering operation commands, and whether the purely decoupled nature of the data models requires an intermediary event-messaging system to notify the visualizers of rapid, deep-nested algorithmic state changes during complex operations like tree balancing.

\clearpage